\chapter{La distribución binomial}\label{ch:g11:binom}

Repita el mismo experimento de sí/no varias veces, de forma independiente, y cuente
los éxitos: el resultado \omterm{def:g11:prob:rv}{distribución} --- el \emph{binomio} --- es
el discreto más importante \omterm{def:g11:prob:rv}{distribución} de todos. Este capítulo lo construye.
con \omterm{met:g10:proba:tree}{árboles} y conteo de caminos; la fórmula cerrada para el camino cuenta (con
factoriales) viene con las herramientas de conteo de \cref{ch:g12:comb}, y el
\omterm{def:g11:prob:rv}{distribución} es revisado en \cref{ch:g12:randvar}.

\section{pruebas de Bernoulli}

\begin{definition}[prueba de Bernoulli]\label{def:g11:binom:bernoulli}
A \emph{prueba de Bernoulli}\index{prueba de Bernoulli} es un experimento con
exactamente dos resultados: \emph{éxito}, con \omterm{def:g10:proba:distribution}{probabilidad} $p $, y
\emph{falla}, con \omterm{def:g10:proba:distribution}{probabilidad} $1 - p $. El variable aleatoria $ X $ igual
a $1$ sobre el éxito y $0$ en caso de fracaso se dice que sigue el
\emph{bernoulli \omterm{def:g11:prob:rv}{distribución}} $\mathcal B(p)$; entonces
\[
\E(X) = p,
\qquad
\V(X) = p(1 - p).
\]
\end{definition}

\begin{proof}[Prueba de las dos fórmulas]
$\E(X) = p \times 1 + (1-p) \times 0 = p $; y desde $ X^2 = X $ (ambos $0$
y $1$ son sus propios cuadrados), $\E(X^2) = p $, entonces por
\cref{prop:g11:prob:konig}, $\V(X) = p - p^2 = p(1-p)$.
\end{proof}

\begin{definition}[Ensayos independientes repetidos]\label{def:g11:binom:repeated}
repitiendo un \omterm{def:g11:binom:bernoulli}{prueba de Bernoulli} $ n $ veces \emph{independientemente} significa: el
El resultado de cada ensayo no tiene influencia sobre los demás, y el \omterm{def:g10:proba:distribution}{probabilidad}
de cualquier completo \omterm{def:g11:seq:sequence}{sucesión} de resultados es el \emph{producto} de la
\omterm{def:g10:proba:distribution}{probabilidades} a lo largo de la trayectoria correspondiente de la \omterm{met:g10:proba:tree}{árbol} --- $ p $ para cada uno
éxito, $1 - p $ por cada fracaso.
\end{definition}

\begin{example}\label{ex:g11:binom:threetrials}
Tres ensayos independientes con éxito \omterm{def:g10:proba:distribution}{probabilidad} $ p $. La \omterm{def:g11:seq:sequence}{sucesión} SFS
(éxito, fracaso, éxito) tiene \omterm{def:g10:proba:distribution}{probabilidad} $ p(1-p)p = p^2(1-p)$ --- y
también lo hace \emph{cada} \omterm{def:g11:seq:sequence}{sucesión} con exactamente dos éxitos, independientemente de
las posiciones: sólo el \emph{número} de los asuntos de S y F.
\end{example}

\section{Recuentos de caminos y coeficientes binomiales.}

\begin{definition}[Coeficiente binomial]\label{def:g11:binom:coefficient}
En el \omterm{met:g10:proba:tree}{árbol} de $ n $ ensayos independientes, el \emph{binomio
coeficiente}\index{binomio
coeficiente} $\binom{n}{k}$ (leer ``$ n $
elegir $ k $'') es el número de rutas que contienen exactamente $ k $ éxitos.
\end{definition}

\begin{example}
$\binom{3}{2} = 3$: las rutas SSF, SFS, FSS. Asimismo $\binom{3}{0} = 1$
(el camino FFF), $\binom{3}{1} = 3$ y $\binom{3}{3} = 1$. Por
convención y por la \omterm{met:g10:proba:tree}{árbol}, $\binom n0 = \binom nn = 1$ por cada $ n $.
\end{example}

\begin{omfigure}
\begin{tikzpicture}[font=\small]
  \coordinate (root) at (0,0);
  \node (S) at (1.5,1.5) {S};
  \node (F) at (1.5,-1.5) {F};
  \node (SS) at (3.0,2.25) {S};
  \node (SF) at (3.0,0.75) {F};
  \node (FS) at (3.0,-0.75) {S};
  \node (FF) at (3.0,-2.25) {F};
  \node (SSS) at (4.5,2.6) {S};
  \node (SSF) at (4.5,1.85) {F};
  \node (SFS) at (4.5,1.1) {S};
  \node (SFF) at (4.5,0.35) {F};
  \node (FSS) at (4.5,-0.35) {S};
  \node (FSF) at (4.5,-1.1) {F};
  \node (FFS) at (4.5,-1.85) {S};
  \node (FFF) at (4.5,-2.6) {F};
  \draw[omThm, thick] (root) -- (S);
  \draw[omThm, thick] (S) -- (SS);   \draw[omThm, thick] (SS) -- (SSF);
  \draw[omThm, thick] (S) -- (SF);   \draw[omThm, thick] (SF) -- (SFS);
  \draw[omThm, thick] (root) -- (F);
  \draw[omThm, thick] (F) -- (FS);   \draw[omThm, thick] (FS) -- (FSS);
  \draw[gray] (SS) -- (SSS);
  \draw[gray] (SF) -- (SFF);
  \draw[gray] (FS) -- (FSF);
  \draw[gray] (F) -- (FF);
  \draw[gray] (FF) -- (FFS);
  \draw[gray] (FF) -- (FFF);
  \node[gray, anchor=west] at (5.0,2.6) {$3$ successes};
  \node[omThm, anchor=west] at (5.0,1.85) {$2$ successes};
  \node[omThm, anchor=west] at (5.0,1.1) {$2$ successes};
  \node[gray, anchor=west] at (5.0,0.35) {$1$ success};
  \node[omThm, anchor=west] at (5.0,-0.35) {$2$ successes};
  \node[gray, anchor=west] at (5.0,-1.1) {$1$ success};
  \node[gray, anchor=west] at (5.0,-1.85) {$1$ success};
  \node[gray, anchor=west] at (5.0,-2.6) {$0$ successes};
\end{tikzpicture}

{\small El \omterm{met:g10:proba:tree}{árbol} de $ n = 3$ ensayos: $\binom{3}{2} = 3$ caminos (rojo) llevan
exactamente dos éxitos, cada uno con \omterm{def:g10:proba:distribution}{probabilidad} $ p^2(1-p)$.}
\end{omfigure}

\begin{proposition}[regla de pascal]\label{prop:g11:binom:pascal}
\index{el triangulo de pascal}
Para $1 \leq k \leq n - 1$:
\[
\binom{n}{k} = \binom{n-1}{k-1} + \binom{n-1}{k}.
\]
\end{proposition}

\begin{proof}
Ordenar los caminos de la $ n $-ensayo \omterm{met:g10:proba:tree}{árbol} con $ k $ éxitos según
su \emph{último} ensayo. Los que terminan en éxito se obtienen de una
camino de la primera $ n-1$ ensayos con $ k - 1$ éxitos: hay
$\binom{n-1}{k-1}$ de ellos. Aquellos que terminan en un fracaso extienden un camino con
$ k $ éxitos entre los primeros $ n-1$ ensayos: $\binom{n-1}{k}$ de ellos.
Cada camino es exactamente de uno de los dos tipos.
\end{proof}

La regla de Pascal genera los coeficientes fila por fila --- cada entrada es el
suma de los dos anteriores:
\[
\begin{array}{ccccccccccc}
&&&&&1&&&&&\\
&&&&1&&1&&&&\\
&&&1&&2&&1&&&\\
&&1&&3&&3&&1&&\\
&1&&4&&6&&4&&1&\\
1&&5&&10&&10&&5&&1
\end{array}
\]

\begin{remark}
Una fórmula cerrada, $\binom nk = \frac{n!}{k!(n-k)!}$, junto con un
teoría sistemática del conteo, se establece en \cref{ch:g12:comb}. en
este nivel, \omterm{prop:g11:binom:pascal}{el triangulo de pascal} calcula todos los coeficientes que necesitamos.
\end{remark}

\section{La distribución binomial}

\begin{theorem}[Distribución binomial]\label{thm:g11:binom:binomial}
Sea $ X $ contar los éxitos en $ n $ independiente pruebas de Bernoulli de
parámetro $ p $. Entonces $ X $ sigue el \emph{binomio
distribución}\index{binomio
distribución} $\mathcal B(n, p)$:
\[
\P(X = k) = \binom{n}{k}\, p^k (1-p)^{n-k},
\qquad k = 0, 1, \dots, n .
\]
\end{theorem}

\begin{proof}
El suceso $ X = k $ es la colección de todos los caminos con exactamente $ k $
éxitos. Cada uno de estos caminos tiene \omterm{def:g10:proba:distribution}{probabilidad} $ p^k(1-p)^{n-k}$: el producto
a lo largo del camino contiene $ k $ factores $ p $ y $ n - k $ factores $1-p $, en
algún orden (\cref{def:g11:binom:repeated}). Hay $\binom nk $ tal
caminos (\cref{def:g11:binom:coefficient}), y sus \omterm{def:g10:proba:distribution}{probabilidades} se suman.
\end{proof}

\begin{example}\label{ex:g11:binom:quiz}
un cuestionario tiene $5$ preguntas independientes, cada una con $4$ opciones; un estudiante
respuestas al azar, por lo que cada pregunta es un éxito con $ p = \frac14$. el
numero $ X $ de respuestas correctas sigue $\mathcal B\left(5, \frac14\right)$,
y, usando fila $5$ de \omterm{prop:g11:binom:pascal}{el triangulo de pascal}:
\[
\P(X = 2) = \binom52 \left(\frac14\right)^2\left(\frac34\right)^3
= 10 \times \frac{1}{16} \times \frac{27}{64}
= \frac{270}{1024} \approx 0.26 .
\]
El \omterm{def:g10:proba:distribution}{probabilidad} de al menos una respuesta correcta utiliza el \omterm{def:g10:proba:operations}{complemento}:
$\P(X \geq 1) = 1 - \P(X = 0) = 1 - \left(\frac34\right)^5 \approx
0.76$.
\end{example}

\begin{proposition}[Esperanza y variación]\label{prop:g11:binom:expectation}
Si $ X \sim \mathcal B(n, p)$:
\[
\E(X) = np,
\qquad
\V(X) = np(1-p).
\]
\end{proposition}

\begin{proof}[Justificación]
escribir $ X = X_1 + X_2 + \dots + X_n $, dónde $ X_i $ es igual $1$ si el
$ i $-ésimo intento tiene éxito: cada uno $ X_i $ es una variable de Bernoulli de \omterm{def:g11:prob:expectation}{esperanza}
$ p $ (\cref{def:g11:binom:bernoulli}). Los promedios suman --- sumando los $ n $
contribuciones da $\E(X) = np $. Eso \emph{variaciones} agregar también para
variables independientes es cierto pero más delicado: la varianza la fórmula es
\emph{admitido en este nivel} y demostrado en \cref{ch:g12:sums}.
\end{proof}

\begin{omfigure}
\begin{tikzpicture}
\begin{axis}[
  width=10cm, height=5cm,
  ybar, bar width=9pt,
  xlabel={$ k $}, ylabel={$\P(X = k)$},
  xmin=-0.8, xmax=10.8, ymin=0, ymax=0.3,
  xtick={0,1,2,3,4,5,6,7,8,9,10},
  ytick={0.1, 0.2},
  tick label style={font=\small},
  label style={font=\small},
  title={$\mathcal B(10,\ 0.5)$}]
  \addplot[fill=omDef!40, draw=omDef] coordinates {
    (0,0.00098) (1,0.00977) (2,0.04395) (3,0.11719) (4,0.20508)
    (5,0.24609) (6,0.20508) (7,0.11719) (8,0.04395) (9,0.00977)
    (10,0.00098)};
\end{axis}
\end{tikzpicture}

{\small El \omterm{def:g11:prob:rv}{distribución} $\mathcal B(10, 0.5)$ (diez lanzamientos de moneda justos):
centrado en $\E(X) = np = 5$, simétrico, con casi todos los \omterm{def:g10:proba:distribution}{probabilidad}
entre $2$ y $8$.}
\end{omfigure}

\begin{method}[Reconocer una situación binomial]\label{met:g11:binom:recognize}
antes de escribir $ X \sim \mathcal B(n, p)$, comprueba tres ingredientes: a
\emph{numero fijo} $ n $ de juicios, decididos de antemano; cada prueba tiene
\emph{dos resultados} con el \emph{mismo} éxito \omterm{def:g10:proba:distribution}{probabilidad} $ p $; el
los ensayos son \emph{independiente} (con repuestos o dispositivos separados).
Extraer sin reemplazo de una población pequeña es \emph{no}
binomio --- el \omterm{def:g10:proba:distribution}{probabilidad} cambios en cada sorteo
(\cref{exo:g11:prob:6}).
\end{method}

\section{Muestreo: ¿es sorprendente la observación?}

El binomio \omterm{def:g11:prob:rv}{distribución} responde una pregunta muy práctica: \emph{si el
éxito \omterm{def:g10:proba:distribution}{probabilidad} realmente lo es $ p $, ¿cuáles son los recuentos de éxitos?
plausible?}

\begin{example}\label{ex:g11:binom:decision}
Se supone que una máquina produce $10\%$ Artículos defectuosos como \omterm{def:g10:functions:extrema}{máximo}. en un
lote de $10$ elementos, $4$ están defectuosos. ¿Mala suerte o máquina rota? si
la máquina está bien, el número de defectos sigue
$\mathcal B(10, 0.1)$, y
\[
\P(X \geq 4) = 1 - \P(X \leq 3) \approx 1 - 0.987 = 0.013 :
\]
aproximadamente una oportunidad en $80$. Observando un suceso esto improbable es un fuerte
señal --- uno \emph{rechaza} la hipótesis de que la máquina todavía funciona
en $10\%$, teniendo en cuenta que la decisión podría ser errónea
\omterm{def:g10:proba:distribution}{probabilidad} acerca de $0.013$.
\end{example}

\begin{method}[Regla de decisión de un modelo binomial]\label{met:g11:binom:decision}
Juzgar un número observado $ k $ de éxitos contra una hipótesis
$ X \sim \mathcal B(n, p)$: calcular el \omterm{def:g10:proba:distribution}{probabilidad}, bajo la hipótesis,
de un resultado \emph{al menos tan extremo} como $ k $. si eso \omterm{def:g10:proba:distribution}{probabilidad} es
muy pequeño (una convención común: a continuación $5\%$), rechazan la hipótesis;
de lo contrario la observación es compatible con ella. El umbral es un
elección, no un teorema: las estadísticas cuantifican el riesgo y el usuario
lo acepta.
\end{method}

\section{Ejercicios}

\begin{exercise}[$\star $]\label{exo:g11:binom:1}
Extender \omterm{prop:g11:binom:pascal}{el triangulo de pascal} remar $7$ y dar los valores de
$\binom62$, $\binom{7}{3}$ y $\binom74$.
\end{exercise}

\begin{exercise}[$\star $]\label{exo:g11:binom:2}
Se lanza un dado justo $4$ veces; $ X $ cuenta los seis. justifica eso
$ X \sim \mathcal B\left(4, \frac16\right)$ y calcular $\P(X = 0)$,
$\P(X = 1)$ y $\P(X \geq 2)$.
\end{exercise}

\begin{exercise}[$\star $]\label{exo:g11:binom:3}
¿Cuál de los siguientes es binomial? Justificar.
\begin{enumerate}
  \item Número de cabezas en $20$ lanzamientos de una moneda justa.
  \item Número de ases en $5$ cartas repartidas de una baraja.
  \item Número de días lluviosos la próxima semana, si cada día llueve con
        \omterm{def:g10:proba:distribution}{probabilidad} $0.3$ independientemente.
\end{enumerate}
\end{exercise}

\begin{exercise}[$\star $]\label{exo:g11:binom:4}
$ X \sim \mathcal B(50, 0.2)$. Dar $\E(X)$, $\V(X)$ y $\sigma(X)$.
\end{exercise}

\begin{exercise}[$\star\star $]\label{exo:g11:binom:5}
Un arquero da en el blanco con \omterm{def:g10:proba:distribution}{probabilidad} $0.7$ en cada disparo,
de forma independiente. En $6$ disparos, calcular el \omterm{def:g10:proba:distribution}{probabilidad} de exactamente $4$ golpes,
y de al menos $5$ golpes.
\end{exercise}

\begin{exercise}[$\star\star $]\label{exo:g11:binom:6}
Una prueba de verdadero/falso tiene $8$ preguntas; un estudiante adivina cada respuesta.
cual es el \omterm{def:g10:proba:distribution}{probabilidad} de pasar (al menos $6$ respuestas correctas)?
\end{exercise}

\begin{exercise}[$\star\star $]\label{exo:g11:binom:7}
Cada caja de cereal comprada contiene la figura A o la figura B, con
\omterm{def:g10:proba:distribution}{probabilidad} $\frac12$ cada uno, de forma independiente. Un coleccionista compra $5$ cajas.
Calcular el \omterm{def:g10:proba:distribution}{probabilidad} que el coleccionista obtenga al menos una figura de
cada tipo. (Complementar: todo A o todo B.)
\end{exercise}

\begin{exercise}[$\star\star $]\label{exo:g11:binom:8}
Un jugador de baloncesto anota tiros libres con \omterm{def:g10:proba:distribution}{probabilidad} $ p $,
de forma independiente. Sea $ X \sim \mathcal B(3, p)$ ser el número de puntuaciones en
tres lanzamientos. Expresar $\P(X = 3)$ y $\P(X \geq 1)$ como funciones de $ p $,
y encontrar para cual $ p $ el \omterm{def:g10:proba:distribution}{probabilidad} de anotar los tres iguales
$\frac{27}{64}$.
\end{exercise}

\begin{exercise}[$\star\star $]\label{exo:g11:binom:9}
¿Cuántas veces se debe lanzar una moneda justa para obtener \omterm{def:g10:proba:distribution}{probabilidad} de conseguir
al menos una cabeza para exceder $0.99$? (Complementar, luego intente sucesivamente
valores de $ n $.)
\end{exercise}

\begin{exercise}[$\star\star $]\label{exo:g11:binom:10}
Usando la regla de Pascal (\cref{prop:g11:binom:pascal}) y
$\binom n0 = \binom nn = 1$, demuestre que las entradas de cada fila de
\omterm{prop:g11:binom:pascal}{el triangulo de pascal} suma a $2^n $: interprete que ambos lados cuentan todos los
caminos de la \omterm{met:g10:proba:tree}{árbol}.
\end{exercise}

\begin{exercise}[$\star\star\star $]\label{exo:g11:binom:11}
Un político afirma $60\%$ aprobación. de forma aleatoria \omterm{def:g10:proba:sample}{muestra} de $10$ gente,
solo $3$ aprobar.
\begin{enumerate}
  \item Según la reclamación, ¿qué \omterm{def:g11:prob:rv}{distribución} hace el numero $ X $ de
        aprobaciones en el \omterm{def:g10:proba:sample}{muestra} ¿seguir? Calcular $\P(X \leq 3)$.
  \item Usando la regla de decisión de \cref{met:g11:binom:decision} con un
        $5\%$ umbral, ¿es la observación compatible con la afirmación?
\end{enumerate}
\end{exercise}

\section{Problema: el tablero Galton}

\begin{problem}[{Problema de fin de semana --- bolas, clavijas y Pascal
Triángulo: cómo nace la forma de campana, por qué favorecen las series de playoffs
el equipo más fuerte, y cuando llorar}]
\label{pb:g11:binom:1}
Deja caer mil bolas a través de un entramado de clavijas, cada una rebota un
Lanzamiento de moneda justo hacia la izquierda o hacia la derecha, y las ranuras de abajo se llenan en un
campana suave y simétrica --- cada vez. La máquina se llama
El tablero de Galton y sus matemáticas son exactamente el tema de este capítulo.
binomio \omterm{def:g11:prob:rv}{distribución} (\cref{thm:g11:binom:binomial}). esto
problema construye el triángulo, recorre el tablero, arbitra un
serie al mejor de siete, y termina donde el binomio gana su
salario: decidir cuando una observación debe hacernos dudar de una
reclamo.

\medskip
\noindent\textbf{Part I --- The triangle.}
\begin{enumerate}
  \item Construir \omterm{prop:g11:binom:pascal}{el triangulo de pascal} abajo a remar $6$
        (\cref{prop:g11:binom:pascal}). Enuncie y explique el
        simetría $\binom nk = \binom{n}{n-k}$ en una frase
        (eligiendo $ k $ objetos es lo mismo que\,\dots).
  \item Verificar en filas $4$ y $5$ que cada fila suma $2^n $,
        y pruébalo: ¿qué hacen todos los $\binom nk $ juntos
        contar?
  \item Volver a derivar la regla de Pascal
        $\binom{n+1}{k} = \binom nk + \binom{n}{k-1}$ por el
        argumento del comité: arreglar a una persona distinguida y
        dividir los comités según el destino de esa persona.
  \item Calcular $\binom73$ dos veces: desde el triángulo, y desde
        la fórmula factorial.
  \item Comprueba la identidad de la escalera.
        $\binom22 + \binom32 + \binom42 + \binom52 =
        \binom63$ y explíquelo poniendo en cascada la regla de Pascal
        de $\binom63$ hacia abajo.
\end{enumerate}

\medskip
\noindent\textbf{Part II --- The board.}
Una pelota cae $ n $ filas de clavijas; en cada clavija rebota
izquierda o derecha con \omterm{def:g10:proba:distribution}{probabilidad} $\frac12$, de forma independiente. Número
las tragamonedas $0$ a $ n $ por el conteo de rebotes derechos.
\begin{enumerate}[resume]
  \item Explique, con la lista de verificación de
        \cref{met:g11:binom:recognize}, por qué el número de ranura
        sigue el binomio \omterm{def:g11:prob:rv}{distribución}
        $\mathcal B\!\left(n, \frac12\right)$.
  \item Para una tabla pequeña ($ n = 4$): dale el quinto espacio
        \omterm{def:g10:proba:distribution}{probabilidades}. ¿Qué franja horaria está más concurrida?
  \item Ahora $ n = 10$ y $1\,024$ bolas: ¿cuáles son las
        \emph{esperado} cuenta de bolas en la ranura del medio, en la ranura
        $7$, y en cada ranura de borde? Describe la forma de la pila.
  \item Para $ X \sim \mathcal B\!\left(10, \frac12\right)$:
        calcular $\E(X)$, $ V(X)$ y $\sigma $
        (\cref{prop:g11:binom:expectation}); luego calcula el
        proporción de bolas esperadas dentro $2\sigma $ de la
        centro (ranuras $2$ a $8$) y comparar con el de Chebyshev
        garantía de \cref{pb:g11:stat:1}.
  \item Una tabla inclinada rebota hacia la derecha con \omterm{def:g10:proba:distribution}{probabilidad} $0.6$:
        dar $\E $, $ V $ y $\sigma $ para $ n = 10$y describir
        ¿Qué pasa con la pila?
  \item En una o dos frases: qué, en el diseño del tablero,
        fabrica la forma de campana --- y ¿por qué tantos?
        cantidades del mundo real (alturas, errores de medición)
        ¿Se acumulan de la misma manera? (El teorema profundo detrás de ambos es
        el teorema del límite central, la cumbre de la universidad
        volúmenes' \omterm{def:g10:proba:distribution}{probabilidad} curso.)
\end{enumerate}

\medskip
\noindent\textbf{Part III --- Best of seven.}
Dos equipos juegan una serie: el primero en $4$ gana y se lleva el título;
Los juegos son independientes.
\begin{enumerate}[resume]
  \item Equipos iguales ($ p = \frac12$): calcula el \omterm{def:g10:proba:distribution}{probabilidad}
        que la serie termina arrasada (exactamente $4$ juegos).
  \item Calcular el \omterm{def:g10:proba:distribution}{probabilidad} que la serie vaya completa
        $7$ juegos (¿cuál debe ser la puntuación después? $6$?).
  \item Completa el \omterm{def:g11:prob:rv}{distribución} de la duración de la serie
        ($4$, $5$, $6$ o $7$ juegos) para equipos iguales, y
        calcular la longitud esperada. ¿Qué longitudes son más
        ¿probablemente?
  \item Ahora un equipo gana cada juego con $ p = 0.6$. Calcular su
        \omterm{def:g10:proba:distribution}{probabilidad} de llevarse la serie (ganar en $4$, $5$, $6$
        o $7$: en cada caso el equipo gana el último juego y
        $3$ de los anteriores). ¿Qué le hizo la serie?
        ¿La ventaja por juego?
  \item Comparar con una sola final ($60\,\%$) y al mejor de 3
        (calcularlo). Indique el efecto general de la longitud de la serie.
        sobre la habilidad contra la suerte --- y por qué las ligas prefieren las largas
        finales.
\end{enumerate}

\medskip
\noindent\textbf{Part IV --- When to cry foul.}
\begin{enumerate}[resume]
  \item Se lanza una moneda $100$ tiempos y espectáculos $62$ cabezas. Para
        una moneda justa, dale $\E $, $\sigma $, y el \omterm{pb:g11:stat:1}{puntuación z}
        (\cref{pb:g11:stat:1}) de la observación. Veredicto
        bajo el $2\sigma $ ¿convención?
  \item Un proveedor afirma $2\,\%$ piezas defectuosas. en un lote
        de $50$ encuentras $3$ defectuoso. Calcular
        $\P(X \geq 3)$ bajo el reclamo
        ($ X \sim \mathcal B(50, 0.02)$; pasar a través
        $\P(X = 0), \P(X = 1), \P(X = 2)$). alarmante ante el
        $5\,\%$ límite (\cref{met:g11:binom:decision},
        \cref{exo:g11:binom:11})?
  \item Persistencia de la lotería: cada billete gana (algo) con
        \omterm{def:g10:proba:distribution}{probabilidad} $\frac{1}{1000}$. Calcular el \omterm{def:g10:proba:distribution}{probabilidad}
        de al menos una victoria con $1\,000$ entradas. la respuesta
        ($\approx 63\,\%$, no $100\,\%$!) esconde un famoso
        constante: calcular $0.999^{1000}$ y mantén el número
        $0.368$ en mente para el grado 12.
  \item Final --- el retrato del binomio: el reconocimiento
        lista de verificación (fija $ n $, independencia, constante $ p $);
        \omterm{prop:g11:binom:pascal}{el triangulo de pascal} como su mesa; la campana como su forma;
        $ np $ y $ np(1-p)$ como su brújula; y sus dos herederos
        esperando en el grado 12 --- la suave curva de la campana y el
        ley de los grandes números. Una frase cada uno.
\end{enumerate}
\end{problem}
